% 中文学术报告 Beamer 模板 —— 组会 / 会议报告
% 引擎：xelatex
\documentclass[UTF8,10pt]{beamer}
\usepackage[fontset=fandol]{ctex}
\usepackage{amsmath,amssymb,amsthm}

\usetheme{Madrid}
\usecolortheme{default}
\setbeamertemplate{navigation symbols}{}
\setbeamertemplate{itemize items}[circle]

% beamer 已内置 theorem/lemma/definition 环境，直接用 block 承载中文定理名
\newcommand{\R}{\mathbb{R}}
\newcommand{\Om}{\Omega}
\DeclareMathOperator{\divg}{div}

\title[简报标题]{非线性偏微分方程解的正则性理论}
\subtitle{Recent Progress and Open Problems}
\author{杨建伟}
\institute{华北水利水电大学 \quad 数学与统计学院}
\date{\today}

\begin{document}

\frame{\titlepage}

\begin{frame}{报告提纲}
\tableofcontents
\end{frame}

\section{问题背景}
\begin{frame}{模型方程}
考虑有界区域 $\Om\subset\R^n$ 上的 Dirichlet 问题：
\begin{equation}\label{eq:model}
  -\divg\bigl(A(x,\nabla u)\bigr)=f \quad \text{in }\Om,
  \qquad u|_{\partial\Om}=0 .
\end{equation}

\begin{itemize}
  \item $A$ 满足 $p$-增长与强制性条件 \only<2>{\alert{（结构性假设）}}
  \item $f\in L^q(\Om)$，$q>n/p$
  \item 核心问题：\alert{弱解是否有界、是否 Hölder 连续？}
\end{itemize}
\end{frame}

\section{主要结果}
\begin{frame}{主要定理}
\begin{block}{主要定理}
设 $f\in L^q(\Om)$ 且 $q>n/p$，$u\in W^{1,p}_0(\Om)$ 是
\eqref{eq:model} 的弱解，则
\begin{equation}
  u\in L^\infty_{\mathrm{loc}}(\Om)\cap C^{0,\alpha}_{\mathrm{loc}}(\Om)
  \quad\text{for some } \alpha\in(0,1).
\end{equation}
\end{block}

\pause
\begin{block}{证明思路}
Caccioppoli 不等式 $\longrightarrow$ De Giorgi 水平集递推
$\longrightarrow$ Moser 迭代 $\longrightarrow$ 振荡衰减。
\end{block}
\end{frame}

\begin{frame}{关键估计}
\begin{columns}[T]
\column{0.48\textwidth}
\textbf{Caccioppoli 型不等式}
\begin{align*}
  \int_{B_r}|\nabla u|^p
  &\le \frac{C}{r^p}\int_{B_{2r}}|u-k|^p \\
  &\quad + C\int_{B_{2r}}|f|^{p'} .
\end{align*}
\column{0.48\textwidth}
\textbf{递推关系}
\begin{align*}
  Y_{j+1} &\le C b^j Y_j^{1+\delta},\\
  Y_j &\to 0 \quad (j\to\infty).
\end{align*}
由 $Y_0$ 充分小推出 $\sup u<\infty$。
\end{columns}
\end{frame}

\section{展望}
\begin{frame}{开放问题与后续工作}
\begin{enumerate}
  \item 临界指数 $q=n/p$ 时的反例构造
  \item 各向异性系数矩阵 $A(x,\xi)$ 的情形
  \item 非局部算子（分数阶 Laplacian）的对应理论
\end{enumerate}

\vspace{1em}
\begin{center}
\Large\emph{谢谢！欢迎批评指正}\\[4pt]
\normalsize Thank you for your attention.
\end{center}
\end{frame}

\end{document}
